\documentclass[twocolumn]{aastex631}
\begin{document}
\title{The reionization ionizing-photon-budget shortfall for star-forming galaxies is not robust to systematics at z\~{}6 (highC systematic corner)}
\author{NebulaMind Lab (autonomous pipeline)}
\affiliation{NebulaMind Lab, \url{https://lab.nebulamind.net}}
\begin{abstract}
Reionization ionizing-photon-budget reconciliation at z\~{}6: star-forming galaxies require f\_esc=0.070 (+0.066/-0.035) to close the budget (Madau-Dickinson SFRD, log xi\_ion=25.5+/-0.15, clumping C=2-5, JWST-SFRD tail), versus indirect-proxy-inferred f\_esc=0.062 (+0.108/-0.039) from LzLCS O32/beta calibrations. Median delta(required-inferred)=+0.006 dex-frac (16-84\%: -0.100 to +0.077); 54\% of the systematic MC shows a shortfall. Conclusion: the budget CLOSES within the systematic, and the sign flips sign between the O32 and beta calibrations (calibration-driven, not robust). The result is bounded by the xi\_ion x clumping x proxy-calibration systematic, not by statistics. Generated autonomously from public data (jwst) via the NebulaMind Lab runner: a bounded, reproducible, descriptive study.
\end{abstract}
\section{Introduction}
Recent studies have highlighted a potential crisis in the reionization photon budget, suggesting that current models may not account for the necessary ionizing photons to drive cosmic reionization [Muñoz2024]. This discrepancy has sparked interest in reassessing the contribution of star-forming galaxies to the ionizing photon budget. Previous research has emphasized the importance of accurately calibrating excursion set reionization models to conserve ionizing photons [Park2022] and understanding the role of galaxy ionizing photon budgets at high redshifts [Duncan2015]. Additionally, Madau \& Dickinson's analytic fitting function for cosmic star formation rate density (SFRD) has been widely used to estimate the ionizing photon production rate [Madau2017].
\section{Data and method}
To address this issue, we employ a literature-anchored budget calculation that does not rely on survey catalog data. Instead, we use the Madau \& Dickinson (2014) analytic fitting function for cosmic SFRD and adopt published values for xi\_ion and O32/beta f\_esc proxy calibrations from LzLCS [Chisholm+22, Flury+22; Simmonds+24]. Our method focuses on reconciling the reionization ionizing-photon-budget using these established parameters.
\section{Result}
Our analysis reveals that star-forming galaxies require an escape fraction of f\_esc=0.070 (+0.066/-0.035) to close the budget at z\~{}6, assuming a Madau-Dickinson SFRD, log xi\_ion=25.5+/-0.15, and clumping factor C=2-5. This value is compared to the indirect-proxy-inferred f\_esc=0.062 (+0.108/-0.039) derived from LzLCS O32/beta calibrations. The median difference between the required and inferred escape fractions is +0.006 dex-frac, with a 16-84\% range of -0.100 to +0.077. Notably, 54\% of our systematic Monte Carlo simulations show a shortfall in the photon budget.
\begin{figure}\centering\includegraphics[width=\columnwidth]{result.png}\caption{The reionization ionizing-photon-budget shortfall for star-forming galaxies is not robust to systematics at z\~{}6 (highC systematic corner)}\end{figure}
\section{Caveats}
It is crucial to acknowledge the limitations of our approach. Our analysis relies on automated, single-selection, and uncalibrated measurements, which may introduce biases and uncertainties. The accuracy of our results depends heavily on the adopted values for xi\_ion and O32/beta f\_esc proxy calibrations, as well as the clumping factor C. Furthermore, our study does not account for potential variations in these parameters across different galaxy populations or environments. Therefore, while our findings provide valuable insights into the reionization photon budget, they should be interpreted with caution and considered alongside other observational and theoretical studies to obtain a more comprehensive understanding of cosmic reionization.
\section*{References}
\footnotesize
[Muoz2024] Muñoz et al. (2024). Reionization after JWST: a photon budget crisis?. bibcode 2024MNRAS.535L..37M \par [Davies2021] Davies et al. (2021). The Predicament of Absorption-dominated Reionization: Increased Demands on Ionizing Sources. bibcode 2021ApJ...918L..35D \par [Park2022] Park et al. (2022). Calibrating excursion set reionization models to approximately conserve ionizing photons. bibcode 2022MNRAS.517..192P \par [Duncan2015] Duncan et al. (2015). Powering reionization: assessing the galaxy ionizing photon budget at z < 10. bibcode 2015MNRAS.451.2030D \par [Madau2017] Madau (2017). Cosmic Reionization after Planck and before JWST: An Analytic Approach. bibcode 2017ApJ...851...50M

\end{document}
